\documentclass{beamer}

\input{settings.tex}


\title{Time Delay}
\subtitle{Advanced Control}
\author{by Sergei Savin}
\centering
\date{\mydate}



\begin{document}
\maketitle


%\begin{frame}{Content}
%\begin{itemize}
%\item Critical point (node)
%\item Stability
%\item Asymptotic stability
%\item Stability vs Asymptotic stability
%\item LTI and autonomous LTI
%\item Stability of autonomous LTI
%\item Read more
%\end{itemize}
%
%\end{frame}



\begin{frame}{Autonomous LTI system with delay}
	\begin{flushleft}
		
		Consider an autonomous delayed LTI system:
		
		\begin{equation}
			\dot{\mathbf{x}}(t) = \mathbf{A} \mathbf{x}(t) + \mathbf{A}_d \mathbf{x}(t - \delta)
		\end{equation}
		
		where $\delta > 0$ is a time delay. This is a \emph{delayed differential equation (DDE)}.
		
		\bigskip
		
		We can pose the following questions:
		
		\begin{itemize}
			\item How do we solve it?
			
			\item How do we analyze stability of the systems of this type?
			
			\item How do we design control for these systems?
		\end{itemize}
		
		
		
	\end{flushleft}
\end{frame}



\begin{frame}{Solution - method of steps}
	\begin{flushleft}
		{\small
		Consider a scalar pure delay system:
		
		\begin{equation}
			\dot{x}(t) = a x(t - \delta)
		\end{equation}
		
		
		In order to solve this system (simulate it forward) we can use a \emph{method of steps}. We specify an initial history function  $\phi(t)$ (which acts as initia conditions) defined on the interval $[ -\delta, \  0]$. Note that for an ODE, initial conditions is a number; for DDE, the initial conditions is a function, the "initial history".
		
		\bigskip
		
		We solve the system forward on the interval $[0, \ \delta]$. On the latter interval, the DDE becomes an ODE:
		
		\begin{equation}
			\dot{x}(t) = a \phi(t - \delta),
		\end{equation}
		
		which has an exact solution on the interval:
		
		\begin{equation}
			x(t) = \phi(0) + a \int_{0}^{t} \phi(\tau - \delta) d \tau.
		\end{equation}
		
		Doing the same for other intervals, e.g. $[\delta, \ 2\delta]$ allows us to solve the system forward. More in the appendix.
		}
	\end{flushleft}
\end{frame}


\begin{frame}{Exponential solutions}
	\begin{flushleft}
		
		Let us consider an exponential solution of the scalar DDE:
		
		\begin{equation}
			x(t) = e^{s t}
		\end{equation}
		
		Substituting to the DDE we get:
		
		\begin{equation}
			s e^{s t} = a e^{s (t - \delta)}
		\end{equation}
		
		Since $e^{s (t - \delta)} = e^{st} e^{-s\delta} $ we can simplify:
		
		\begin{equation}
			s  = a e^{-s\delta} 
		\end{equation}
		
		This is a \emph{characteristic equation} of the DDE. %Exponential solutions form a basis in the solution space of this DDE. Any particular solution to the original DDE will be a linear combination of this infinite family of solutions.
		
		Note that $s  = a e^{-s\delta} $ has infinitely many complex roots, each corresponding to a valid solution. 
		
	\end{flushleft}
\end{frame}







\begin{frame}{Stability}
	\begin{flushleft}
		
		A solution of a linear DDE can represented as a weighted sum of its exponential solutions, associated with the roots $s_i$:
		
		\begin{equation}
			x(t) = \sum_{i=1}^\infty c_i e^{s_i t}
		\end{equation}
		
		If at least one of the roots $s_i$ has a positive real part, the system is unstable.
		
		
	\end{flushleft}
\end{frame}





\begin{frame}{LTI System with delay}
\begin{flushleft}

Consider the following LTI with an input delay:

\begin{equation}
    \dot{\mathbf{x}}(t) = \mathbf{A} \mathbf{x}(t) + \mathbf{B} \mathbf{u}(t - \delta)
\end{equation}

where $\delta > 0$ is a constant time delay.

\bigskip

Let us propose a linear control law:

\begin{equation}
	\mathbf{u}(t) = -\mathbf{K} \mathbf{x}(t)
\end{equation}

\bigskip

The closed-loop system is: 

\begin{equation}
	\dot{\mathbf{x}}(t) = \mathbf{A} \mathbf{x}(t) - \mathbf{B} \mathbf{K}  \mathbf{x}(t - \delta)
\end{equation}

We cannot analyse the stability of this system  by simply analysing the eigenvalues of $\mathbf{A} - \mathbf{B}\mathbf{K}$ as we do in the absence of delays.


\end{flushleft}
\end{frame}




\begin{frame}{Exact stability analysis, 1}
	\begin{flushleft}
		
		In order to check stability of the system, we can analyse its exponential solutions:
		%
		\begin{align}
			\mathbf{x}(t) = \mathbf{v} e^{s t} 
		\end{align}
		
		The DDE becomes:
		%
		\begin{align}
			s \mathbf{v} e^{s t}  = \mathbf{A} \mathbf{v} e^{s t}  - \mathbf{B} \mathbf{K}\mathbf{v} e^{s (t - \delta)}
			\\
			s \mathbf{v}  = \mathbf{A} \mathbf{v}  - \mathbf{B} \mathbf{K}\mathbf{v} e^{-s \delta}
			\\
			(s\mathbf{I} - \mathbf{A} + \mathbf{B} \mathbf{K} e^{-s \delta}) \mathbf{v} = 0
		\end{align}
		
		This implies that the matrix on the left has a non-trivial null-space, so:
		%
		\begin{align}
			\text{det}\left(s\mathbf{I} - \mathbf{A} + \mathbf{B} \mathbf{K} e^{-s \delta} \right) = 0
		\end{align}
		
		This is the characteristic equation of the DDE.
		
	\end{flushleft}
\end{frame}



\begin{frame}{Exact stability analysis, 2}
	\begin{flushleft}
		
		The DDE $\dot{\mathbf{x}}(t) = \mathbf{A} \mathbf{x}(t) - \mathbf{B} \mathbf{K}  \mathbf{x}(t - \delta)$ is stable if the characteristic equation:		
		%
		\begin{align}
			\text{det}\left(s\mathbf{I} - \mathbf{A} + \mathbf{B} \mathbf{K} e^{-s \delta} \right) = 0
		\end{align}
		%
		has roots with strictly negative real parts.
		
		
	\end{flushleft}
\end{frame}




\begin{frame}{Small delay approximation}
	\begin{flushleft}
		
		Consider the function $\mathbf{x}(t - \delta)$. We can define $\tau = t - \delta$ and write the Taylor expansion of $\mathbf{x}(\tau)$ around the point $\tau = t$:
		
		\begin{equation}
			\mathbf{x}(\tau) = \mathbf{x}(t) + \frac{\partial \mathbf{x}}{\partial \tau} (\tau - t) + \text{h.o.t}
		\end{equation}
		
		Notice that:
		
		\begin{equation}
			\frac{\partial \mathbf{x}}{\partial t} = 
			\frac{\partial \mathbf{x}}{\partial \tau} \frac{\partial \tau}{\partial t} = 
			\frac{\partial \mathbf{x}}{\partial \tau} 
		\end{equation}
		
		
		Noting that $\tau - t = -\delta$ and dropping the higher-order terms, we get:
		
		\begin{equation}
			\mathbf{x}(t - \delta) \approx \mathbf{x}(t) -\delta \frac{\partial \mathbf{x}}{\partial t}
		\end{equation}
		
		
	\end{flushleft}
\end{frame}



\begin{frame}{Approximated closed-loop system}
	\begin{flushleft}
		
		Substituting the small delay approximation, we approximate the closed-loop system:
		
		\begin{equation}
			\dot{\mathbf{x}}(t) = \mathbf{A} \mathbf{x}(t) - \mathbf{B} \mathbf{K}  
			\textcolor{myblue}{ \left(\mathbf{x}(t) -\delta \dot{\mathbf{x}}(t) \right)}
		\end{equation}
		
		Now, all variables depend on current time only. Rearranging the terms we get:
		
		\begin{align}
			\dot{\mathbf{x}} = \mathbf{A} \mathbf{x} - 
			\mathbf{B} \mathbf{K}  \mathbf{x} + \delta \mathbf{B} \mathbf{K}   \dot{\mathbf{x}}
		\\
			\dot{\mathbf{x}} - \delta \mathbf{B} \mathbf{K}   \dot{\mathbf{x}} = (\mathbf{A} - \mathbf{B} \mathbf{K})  \mathbf{x}
			\\
			\dot{\mathbf{x}} = 
			(\mathbf{I} - \delta \mathbf{B} \mathbf{K})^{-1} (\mathbf{A} - \mathbf{B} \mathbf{K})  \mathbf{x}
		\end{align}
		
		With that, we can analyse the stability of the approximated system via eigenvalues of the matrix $\mathbf{M}  = (\mathbf{I} - \delta \mathbf{B} \mathbf{K})^{-1} (\mathbf{A} - \mathbf{B} \mathbf{K}) $. However, this does not guarantee the stability of the true delyed system.
		
		
	\end{flushleft}
\end{frame}





\begin{frame}{Literature}

\begin{itemize}
\item ***
\end{itemize}

\end{frame}



\myqrframe


\begin{frame}{Appendix: Scalar case (general)}
	\begin{flushleft}
		
		
		Previously we considered a pure-delay system $\dot{x}(t) = a_d x(t - \delta)$. Now we can generalize:
		
		\begin{equation}
			\dot{x}(t) = a x(t) + a_d x(t - \delta)
		\end{equation}
		
		As before, we initialize the system with a function $\phi(t)$ defined on the interval $[ -\delta, \  0]$ and solve the system forward on the interval $[0, \ \delta]$:
		
		\begin{equation}
			\dot{x}(t) = a x + a_d \phi(t - \delta),
		\end{equation}
		
		which has an exact solution on the interval:
		
		\begin{equation}
			x(t) = e^{a t} \phi(0) + ... + a_d \int_{0}^{t} e^{a (t - \tau)} \phi(\tau - \delta) d \tau.
		\end{equation}
		
		Again, doing the same for other intervals, e.g. $[\delta, \ 2\delta]$ allows us to solve the system forward.
		
	\end{flushleft}
\end{frame}




\end{document}
